\section{Introduction}
\label{sec:intro}

When you read ``the garden was full of roses,'' you smell them. This involuntary sensory activation is a hallmark of human language comprehension~\citep{barsalou1999perceptual}---but does anything analogous happen inside a large language model? If an LLM processes the word ``thunder,'' does its internal representation shift toward an auditory direction, even when the surrounding text says nothing about sound?

\para{Why this matters.}
The question strikes at the heart of the symbol grounding debate~\citep{harnad1990symbol}. A growing body of work shows that LLMs can predict human sensory similarity judgments~\citep{marjieh2023large}, produce sensory language in creative writing~\citep{hicke2025zero}, and align with vision and audio encoders under explicit prompting~\citep{wang2025words}. Yet a fundamental gap remains: these results demonstrate that LLMs \emph{can} produce sensory associations when asked, but not that they \emph{spontaneously} activate them during ordinary language processing. The distinction matters because spontaneous activation would suggest that LLMs build an implicit sensory model of the world from text alone---a ``simulated sensory world''---rather than merely retrieving associations on demand.

\para{The gap.}
Prior work has studied explicit sensory prompting~\citep{wang2025words}, sensory language generation~\citep{hicke2025zero}, and prompted sensory judgments~\citep{marjieh2023large, lee2025exploring}. However, no study has systematically tested whether sensory representations activate \emph{implicitly} when LLMs process objects with strong sensory associations in narrative context. Moreover, existing studies typically examine one or two modalities; a comprehensive analysis of all five classic senses plus non-traditional senses (\eg interoception) within a single representational framework is lacking.

\para{Our approach.}
We construct 450 stimulus sentences embedding 150 sensory words (25 per modality) in three conditions: \emph{implicit} (narrative context with no sensory verbs, \eg ``The perfume drifted through the open window''), \emph{explicit} (overt sensory description, \eg ``The perfume smelled incredibly strong''), and \emph{control} (neutral framing, \eg ``The perfume was mentioned in the report''). We extract hidden states from \qwen at all 29 layers and train linear probes to decode the dominant sensory modality. \figref{fig:overview} illustrates the experimental pipeline.

\para{Key results.}
Linear probes decode sensory modality at \textbf{90.0\%} accuracy in the implicit condition---$5.4\times$ above the 16.7\% chance level---and at 92.7\% in the explicit condition, with permutation-test significance at $p = 0.005$. Implicit and explicit representations are highly similar (cosine similarity $= 0.822$), significantly more so than implicit and control representations ($p = 6.1 \times 10^{-20}$). The six senses form mutually opposing directions in representation space (mean pairwise cosine $= -0.20$), and interoception participates in the same geometric structure as the classic five senses.

\para{Contributions.} We make the following contributions:
\begin{itemize}[leftmargin=*,itemsep=0pt,topsep=0pt]
    \item We demonstrate that LLMs spontaneously activate sensory-specific representations when processing sensory-laden objects, even without explicit sensory language---achieving 90\% linear probe accuracy in the implicit condition.
    \item We show that implicit and explicit sensory mentions activate overlapping internal representations, suggesting a persistent sensory encoding rather than post-hoc retrieval.
    \item We characterize the geometry of the sensory subspace, revealing that six modalities (including interoception) form a structured simplex of mutually opposing directions.
    \item We release our stimulus set, probing pipeline, and all experimental code to support reproducibility and future work on sensory representations in language models.
\end{itemize}
